% Artigo final do Lab03 no template SBC.
% Para compilar no Overleaf: crie um projeto a partir do "SBC Conferences
% Template", substitua o main.tex por este arquivo e envie as figuras de
% analise/saidas/ (rq01_04_distribuicoes.png e rq05_dispersao.png).
% Limite do enunciado: 10 páginas. Confira a contagem após compilar.
\documentclass[12pt]{article}

\usepackage{sbc-template}
\usepackage{graphicx,url}
\usepackage[utf8]{inputenc}
\usepackage[brazil]{babel}
\usepackage{booktabs}

\sloppy

\title{Métricas DORA em repositórios populares do GitHub:\\
quanto as conclusões dependem da definição operacional?}

% TODO(dupla): preencher nomes e e-mails.
\author{Integrante A\inst{1}, Integrante B\inst{1}}

\address{Pontifícia Universidade Católica de Minas Gerais (PUC Minas)\\
  Laboratório de Experimentação de Software\\
  Belo Horizonte -- MG -- Brasil
}

\begin{document}

\maketitle

\begin{abstract}
We mined 306 popular GitHub repositories that use GitHub Actions and computed
approximations of the four DORA metrics over a 12-month window. The median
repository publishes one release every three weeks, has a lead time of about
two weeks and recovers from a CI failure in under four hours. A gold sample
of 300 releases labelled by two raters ($\kappa=0.85$) showed that a
version/commit-based heuristic for corrective releases has precision 0.38;
a revised heuristic based on release notes reaches F1 0.59 on held-out
repositories. Swapping operational definitions changes the overall DORA
category of 21--38\% of the repositories.
\end{abstract}

\begin{resumo}
Mineramos 306 repositórios populares do GitHub que usam GitHub Actions e
calculamos aproximações das quatro métricas DORA em uma janela de 12 meses.
O repositório mediano publica uma release a cada três semanas, tem lead time
de cerca de duas semanas e se recupera de uma falha de CI em menos de quatro
horas. Uma amostra-ouro de 300 releases rotuladas por dois avaliadores
($\kappa=0{,}85$) mostrou que uma heurística de release corretiva baseada em
versão e commits tem precisão de 0,38; uma versão revista, baseada nas notas
de release, chega a F1 de 0,59 em repositórios separados para teste. Trocar
as definições operacionais muda a categoria DORA geral de 21\% a 38\% dos
repositórios.
\end{resumo}

\section{Introdução}

As métricas DORA --- frequência de deploy, lead time, taxa de falha de
mudanças (CFR) e tempo de recuperação --- são usadas para comparar o
desempenho de entrega de software. Em repositórios públicos, nenhuma delas é
observada diretamente: é preciso escolher \emph{proxies}, como tratar uma
release publicada como deploy ou um run de CI vermelho como falha. Este
trabalho mede as métricas em 306 repositórios e avalia quanto as conclusões
dependem dessas escolhas. As hipóteses abaixo foram registradas antes da
análise.

\textbf{RQ01 (frequência):} mediana na faixa Medium, com cauda longa.
\textbf{RQ02 (lead time):} mediana entre uma semana e um mês; a variante por
release maior que a por commit na maioria dos repositórios.
\textbf{RQ03 (CFR):} o proxy de CI maior que o de releases corretivas.
\textbf{RQ04 (recuperação):} mediana abaixo de 24 horas, com alguma censura.
\textbf{RQ05 (frequência $\times$ CFR):} correlação fraca com o CFR de CI e
positiva com o de releases, por efeito mecânico.
\textbf{RQ06 (fatores):} diferenças entre linguagens, popularidade, idade e
tipo de projeto.
\textbf{RQ07 (sensibilidade):} ao menos uma troca de definição muda a
categoria de parte dos repositórios.
\textbf{RQ08 (rework rate):} maior que o CFR por releases e alto em termos
absolutos, por causa da heurística.

\section{Metodologia}

\textbf{Fonte e janela.} Os dados vêm da API REST do GitHub, por um cliente
próprio com cache em disco, paginação e espera por cota. A janela vai de
1/10/2025 a 30/9/2026. A coleta foi feita em outubro de 2026.

\textbf{Seleção.} A busca por repositórios com pelo menos 1.000 estrelas, não
arquivados e com \emph{push} na janela foi fatiada em oito faixas de estrelas,
porque cada consulta devolve no máximo 1.000 resultados. Os candidatos foram
embaralhados com semente 42 e avaliados até a amostra chegar a 306. Cada
candidato precisava ter GitHub Actions, pelo menos 5 releases publicadas (sem
draft nem pré-release) e 50 workflow runs válidos no \emph{default branch} na
janela. A Tabela~\ref{tab:funil} resume o funil.

\begin{table}[ht]
\centering
\caption{Funil de seleção.}\label{tab:funil}
\small
\begin{tabular}{lrrl}
\toprule
Etapa & Restam & Saem & Motivo principal \\
\midrule
Encontrados na busca & 37.401 & -- & soma das oito faixas \\
Candidatos acessíveis & 7.451 & 29.950 & teto de 1.000 por consulta \\
Avaliados & 1.128 & 6.323 & amostra completada antes \\
Com GitHub Actions & 964 & 164 & sem workflows \\
$\geq$5 releases e $\geq$50 runs & \textbf{306} & 658 & 570 com menos de 5 releases \\
\bottomrule
\end{tabular}
\end{table}

\textbf{Definições.} Na combinação de referência (C1), deploy é uma release
principal; a frequência é o número delas por semana. O lead time (a) é a
mediana, entre as releases, da publicação menos o commit mais antigo do
\texttt{compare}; o (b) é a mediana sobre todos os commits. O CFR (a) é a
proporção de runs de CI com falha. O CFR (b) é a proporção de releases
seguidas por uma release corretiva em até sete dias, com censura dos últimos
sete dias da janela. O tempo de recuperação é a mediana, por workflow, entre
a primeira falha e o sucesso seguinte. O rework rate (RQ08) é a proporção de
releases corretivas. Para a RQ07, C2 usa releases e pré-releases, lead time
(b) e CFR (b); C3 usa tags datadas pelo commit apontado, lead time (b) e
CFR (a).

\textbf{Validação manual.} Foram sorteados 60 repositórios e 5 releases de
cada (semente 42). Dois avaliadores rotularam de forma independente, com
protocolo escrito: tipo do projeto, se as releases são entregas reais e se
cada release é corretiva (objetivo principal de corrigir defeito, julgado
pelas notas e commits, não pelo número da versão). A concordância (kappa de
Cohen) foi substancial para tipo ($\kappa=0{,}69$), moderada para entregas
reais ($\kappa=0{,}41$) e quase perfeita para corretividade
($\kappa=0{,}85$). As 43 divergências (13, 9 e 21) foram decididas
item a item com apoio de um assistente de IA, que aplicou os critérios do
protocolo às evidências; a justificativa de cada decisão está no
repositório. O consenso tem 107 releases corretivas (35,7\%).

\textbf{Heurística corretiva.} A v1 marca uma release como corretiva se o
SemVer muda só o patch ou se algum commit contém \emph{fix}, \emph{hotfix},
\emph{revert} ou \emph{bugfix}. Para revisá-la sem ajustar aos itens de
avaliação, os 60 repositórios foram divididos em metades de desenvolvimento e
teste; a regra foi escolhida na primeira e medida uma única vez na segunda.
A v2 marca a release quando as notas são dominadas por correções (seção de
correções sem seção de funcionalidades, ou mais itens de correção que de
funcionalidade) ou quando é um patch sem funcionalidade nas notas e com
commits de correção em número igual ou maior que os de funcionalidade
(Tabela~\ref{tab:heur}). A v2 é usada no dataset final, mas não atinge a
meta de F1 $\geq 0{,}70$.

\begin{table}[ht]
\centering
\caption{Heurística corretiva contra o consenso.}\label{tab:heur}
\small
\begin{tabular}{llrrrr}
\toprule
Versão & Conjunto & $n$ & Precisão & Recall & F1 \\
\midrule
v1 & todos & 300 & 0,375 & 0,944 & 0,537 \\
v1 & teste & 150 & 0,290 & 0,950 & 0,444 \\
v2 & desenvolvimento & 150 & 0,624 & 0,791 & 0,697 \\
v2 & \textbf{teste} & \textbf{150} & \textbf{0,500} & \textbf{0,725} & \textbf{0,592} \\
\bottomrule
\end{tabular}
\end{table}

\textbf{Análise.} Mediana e IQR por métrica; Spearman na RQ05;
Kruskal-Wallis com correção de Holm sobre os 24 testes, $\varepsilon^2$ e
Cliff's $\delta$ (descritivo) na RQ06; percentual de mudança de categoria e
kappa ponderado linear na RQ07.

\section{Resultados e discussão}

\begin{table}[ht]
\centering
\caption{Métricas da amostra (mediana e quartis).}\label{tab:metricas}
\small
\begin{tabular}{lrrrr}
\toprule
Métrica & $n$ & Mediana & Q1 & Q3 \\
\midrule
Frequência (releases/semana) & 306 & 0,345 & 0,192 & 0,743 \\
Lead time (a) por release (dias) & 306 & 13,98 & 3,01 & 33,94 \\
Lead time (b) por commit (dias) & 306 & 12,24 & 3,44 & 47,52 \\
CFR (a) de CI & 306 & 0,065 & 0,023 & 0,122 \\
CFR (b) de releases & 305 & 0,250 & 0,091 & 0,540 \\
Recuperação (horas) & 284 & 3,75 & 0,93 & 14,95 \\
Rework rate & 306 & 0,542 & 0,344 & 0,729 \\
\bottomrule
\end{tabular}
\end{table}

\textbf{RQ01.} A mediana é de 0,345 release por semana, cerca de uma a cada
três semanas; só 1,3\% dos repositórios chegam a Elite, 46,1\% ficam em
Medium e 33,7\% em Low (Figura~\ref{fig:dist}). A hipótese se confirma: a
release do GitHub é um evento planejado, mais raro que o deploy contínuo.

\textbf{RQ02.} As medianas são 14,0 dias (a) e 12,2 dias (b), ambas em
Medium. A variante (a) é maior em 59,8\% dos repositórios, o que confirma a
hipótese só em parte. A (a) explode com um único commit antigo
(\emph{cherry-pick}); a (b) é dominada pelas releases grandes, que acumulam
muitos commits.

\textbf{RQ03.} O CFR (a) tem mediana de 6,5\%, e 80,4\% dos repositórios
ficam em Elite. O CFR (b) tem mediana de 25,0\%, quase quatro vezes maior; a
hipótese não se confirma. Os proxies medem coisas diferentes: falha de
pipeline, que raramente chega ao usuário, e correção publicada. Com a v1, o
CFR (b) tinha mediana de 33,3\%; como a v2 ainda tem cerca de um falso
positivo por verdadeiro positivo, o valor deve ser lido como limite superior
aproximado.

\textbf{RQ04.} A mediana é de 3,7 horas, e 84,9\% dos repositórios se
recuperam em menos de um dia. Só 0,7\% dos 15.709 episódios ficaram
censurados, mas eles aparecem em 22\% dos repositórios.

\begin{figure}[ht]
\centering
\includegraphics[width=.9\textwidth]{rq01_04_distribuicoes.png}
\caption{Distribuição das métricas DORA na amostra.}\label{fig:dist}
\end{figure}

\textbf{RQ05.} A correlação entre frequência e CFR (a) é fraca
($\rho=0{,}13$, $p=0{,}019$), e com o CFR (b) é forte ($\rho=0{,}68$,
$p<0{,}001$; Figura~\ref{fig:rq05}). A segunda é em boa parte mecânica: mais
releases deixam mais releases a menos de sete dias de uma corretiva. Os dados
não sustentam nem refutam com segurança a afirmação do DORA de que velocidade
e estabilidade andam juntas.

\begin{figure}[ht]
\centering
\includegraphics[width=.75\textwidth]{rq05_dispersao.png}
\caption{Frequência de deploy $\times$ CFR (a) e (b).}\label{fig:rq05}
\end{figure}

\textbf{RQ06.} Após Holm, linguagem se associa a frequência
($\varepsilon^2=0{,}13$), lead time (a) ($0{,}17$) e CFR (b) ($0{,}16$);
quartil de estrelas, só ao lead time (b) ($0{,}06$); e quartil de idade, a
frequência, aos dois lead times ($0{,}22$ e $0{,}23$), ao CFR (b) ($0{,}10$) e
à recuperação ($0{,}08$). O quartil mais jovem publica mais (0,54 contra 0,29
release/semana) e tem CFR (b) mais alto (0,50 contra cerca de 0,20), o mesmo
mecanismo da RQ05. Para tipo do projeto (60 repositórios do consenso), nenhum
teste foi significativo; o menor p ajustado foi o da recuperação (0,059),
com bibliotecas mais lentas (10,5 horas) que aplicações e CLIs (3,2--3,5).
Com grupos pequenos, ausência de significância não indica equivalência.

\textbf{RQ07.} A Tabela~\ref{tab:rq07} mostra a mudança da categoria geral. O
proxy de CFR é a escolha que mais pesa: trocar o CFR (a) pelo (b) muda a
categoria de CFR de 68,9\% dos repositórios. A hipótese se confirma: C2
mudou mais que C3 na categoria geral. A causa, porém, foi o CFR, e não as
pré-releases. Na categoria de frequência, isoladamente, as tags de C3 mudaram
mais repositórios (19,0\%) que as pré-releases de C2 (8,2\%).

\begin{table}[ht]
\centering
\caption{Mudança da categoria DORA geral entre combinações.}\label{tab:rq07}
\small
\begin{tabular}{lrrrr}
\toprule
Par & $n$ & Mudaram & \% & $\kappa$ ponderado \\
\midrule
C1 $\times$ C2 & 283 & 108 & 38,2 & 0,415 \\
C1 $\times$ C3 & 284 & 60 & 21,1 & 0,714 \\
C2 $\times$ C3 & 283 & 100 & 35,3 & 0,500 \\
\bottomrule
\end{tabular}
\end{table}

\textbf{RQ08.} Com a v2, o rework rate tem mediana de 0,54, é maior que o
CFR (b) em 70,2\% dos repositórios (Wilcoxon $p<0{,}001$) e tem correlação
fraca e positiva com ele ($\rho=0{,}30$). Com a v1, a mediana era 0,98 e a
correlação era negativa: a métrica media a presença de algum \emph{fix}, não
o objetivo da release. A hipótese se confirma, e a amostra-ouro mostra que a
v1 não servia para essa métrica.

\section{Ameaças à validade}

\textbf{Construto.} Release não é deploy: a frequência medida é de releases,
e 570 dos descartes do funil foram por falta delas. A amostra-ouro não teve
nenhum repositório com releases rotuladas como entregas não reais (4
incertos). O CFR (a) mede falha de CI, não de produção. O CFR (b) e o rework
dependem de uma heurística com F1 de 0,59. A definição escolhida muda a
categoria de até 38\% dos repositórios (RQ07).

\textbf{Interna.} Idade e volume de releases estão confundidos, e a
associação entre frequência e CFR (b) é em parte mecânica. Os fatores da RQ06
foram testados isoladamente. A v2 depende do estilo das notas de release, que
varia entre comunidades. Episódios censurados e meses no teto de 1.000 runs
(35 repositórios) podem enviesar a recuperação e o CFR (a). As 43
divergências foram decididas com apoio de IA, e não em reunião dos
avaliadores. Elas afetam 10\% dos rótulos, e as justificativas estão
registradas.

\textbf{Externa.} A amostra tem projetos abertos, populares, com Actions e
releases. A busca alcançou só os mais populares de cada faixa (7.451 de 37.401
candidatos). Os resultados descrevem uma janela de 12 meses.

\textbf{Conclusão.} Os 24 testes da RQ06 foram corrigidos por Holm. Várias
linguagens e o tipo \emph{outro} têm menos de cinco repositórios. A metade de
teste da heurística tem 40 releases corretivas, e o F1 estimado tem
incerteza considerável.

\section{Replicação cruzada}

% TODO(dupla): substituir pelo relato real após receber o grupo sorteado.
A replicação do pipeline de outro grupo ainda não foi feita. Falta
identificar o grupo e obter o repositório, o CSV original e a lista de 30
repositórios. O procedimento (execução só pelo README, comparação com
\texttt{replicacao.comparar}, Issues no repositório autor) está documentado
em \texttt{artigo/replicacao\_cruzada.md}. Até 23/10/2026, nenhuma Issue de
grupo replicador havia sido aberta no nosso repositório.

\section{Conclusão}

As métricas DORA podem ser aproximadas com dados públicos do GitHub, mas os
valores e as categorias dependem da definição operacional. O proxy de CFR é a
escolha que mais pesa. A validação manual mostrou que uma heurística ingênua
de release corretiva tem precisão de 0,38. Uma versão baseada nas notas de
release melhora o resultado sem chegar à meta. Comparações de desempenho DORA
entre projetos abertos devem declarar a definição usada e a validação dos
\emph{proxies}.

\section*{Disponibilidade}

Repositório:
\url{https://github.com/gabrieltinoco/Laboratorio-Experimentacao-de-Software}
(pasta \texttt{Laboratorio03}).
% TODO(dupla): link do GitHub Projects.
GitHub Projects: \texttt{<preencher>}.

\bibliographystyle{sbc}
\begin{thebibliography}{9}
\bibitem{cohen} Cohen, J. (1960). A coefficient of agreement for nominal
scales. \emph{Educational and Psychological Measurement}, 20(1):37--46.
\bibitem{dora} DORA (2024). \emph{A history of DORA's software delivery
metrics}. \url{https://dora.dev/insights/dora-metrics-history/}.
\bibitem{landis} Landis, J.~R. and Koch, G.~G. (1977). The measurement of
observer agreement for categorical data. \emph{Biometrics}, 33(1):159--174.
\bibitem{wohlin} Wohlin, C. et al. (2012). \emph{Experimentation in Software
Engineering}. Springer.
\end{thebibliography}

\end{document}
